% !TeX root = ../main.tex
\section{Experimental Evaluation and Results}
\label{sec:evaluacion}

This section quantitatively analyzes the performance of the proposed architecture across the experimental testbed, directly addressing the core research questions: the temporal overhead of the auditability layer, the impact of physical resource decoupling, the latency impact of expanding the IBFT 2.0 validator committee, the empirical resilience against adversarial tampering attacks, and the predictability of EVM gas consumption. To ensure statistical soundness and experimental reproducibility, each of the five experimental configurations ($E_1$ to $E_5$) was evaluated across 3 complete, independent executions of 30 federated rounds each, yielding a corpus of 15 executions and 450 total federated rounds ($N_{\text{total}} = 450$). We note that the model convergence trajectory is identical with or without the blockchain layer, as the traceability framework operates fully decoupled and does not modify training tensors or gradients.

\subsection{Mathematical Formalization of Performance Metrics}
To isolate the latency imposed by each system component, the total round time at the edge nodes ($T_{\text{client}}$) and at the central coordinator ($T_{\text{server}}$) during iteration $t$ is analytically defined as:
\begin{align}
T_{\text{client}}^{(t)} &= t_{\text{verif\_desc}}^{(t)} + t_{\text{train}}^{(t)} + t_{\text{ipfs\_up}}^{(t)} + t_{\text{tx\_besu}}^{(t)} \label{eq:t_client} \\
T_{\text{server}}^{(t)} &= t_{\text{wait}}^{(t)} + t_{\text{verif\_asc}}^{(t)} + t_{\text{fedavg}}^{(t)} \nonumber \\
&\quad + t_{\text{ipfs\_global}}^{(t)} + t_{\text{block\_confirm}}^{(t)} \label{eq:t_server}
\end{align}
where $t_{\text{verif\_desc}}$ denotes local downward verification against the Besu CID; $t_{\text{train}}$, gradient computation time in PyTorch; $t_{\text{ipfs\_up}}$, binary model publication to IPFS; $t_{\text{tx\_besu}}$, transaction issuance and confirmation for \texttt{registerClientModel}; $t_{\text{wait}}$, the coordinator gRPC reception window; $t_{\text{verif\_asc}}$, on-chain validation via \texttt{getClientModel}; $t_{\text{fedavg}}$, matrix aggregation; and $t_{\text{block\_confirm}}$, IBFT 2.0 block confirmation following \texttt{updateModel}.

These two timing perspectives provide complementary system insights:
\begin{itemize}
    \item \textbf{Client cycle time ($T_{\text{client}}$, reported in Table~\ref{tab:tiempos_resumen}):} Quantifies local computation and communication duration at edge nodes (vehicles), from the receipt of the global model to the confirmation of their contribution on the blockchain. This metric governs client operational autonomy and maintains comparability with prior benchmarks~\cite{cloudlabaia_flblockchain}.
    \item \textbf{End-to-end round time at the server ($T_{\text{server}}$, broken down in Table~\ref{tab:desglose_fases}):} Encompasses the entire round orchestrated by the central coordinator, including gRPC client synchronization and block inclusion of the aggregated global model via \texttt{updateModel} ($\approx 2.8\text{--}5.1$~s additional latency per round).
\end{itemize}

Relative temporal overhead ($\Delta_{\text{overhead}}$) between a secure configuration $E_{\text{BC}}$ and its corresponding baseline $E_{\text{base}}$ is calculated using the distribution medians:
\begin{equation}
\Delta_{\text{overhead}} = \frac{\text{Med}(T_{E_{\text{BC}}}) - \text{Med}(T_{E_{\text{base}}})}{\text{Med}(T_{E_{\text{base}}})} \times 100 \quad [\%] \label{eq:overhead}
\end{equation}

Regarding statistical methodology, the experimental protocol treats the \textbf{15 complete executions} as independent sample units ($N_{\text{runs}} = 3$ independent repetitions per configuration). The 30 consecutive rounds of each run are not treated as independent and identically distributed (i.i.d.) observations, because round $t+1$ depends intrinsically on the model state and system conditions at round $t$. Consequently, all 450 rounds (90 per configuration) are utilized for longitudinal descriptive analysis and trajectory visualization (Figure~\ref{fig:tiempo_todas}), while inter-run reproducibility is quantified via the Standard Error ($SE$) across execution means ($SE = \sigma_{\text{runs}} / \sqrt{N_{\text{runs}}}$).

\subsection{Consolidated Temporal Performance Overview}
Table~\ref{tab:tiempos_resumen} summarizes the per-round client cycle time ($T_{\text{client}}$) across the 15 independent runs (3 per configuration, totaling 90 rounds per environment), reporting the median, mean, inter-run standard deviation ($\sigma_{\text{runs}}$), and standard error ($SE = \sigma_{\text{runs}} / \sqrt{3}$). Figure~\ref{fig:tiempo_todas} depicts the continuous time-series trajectories across all rounds.

\begin{table}[htbp]
\caption{Consolidated per-round client cycle time ($T_{\text{client}}$) across 15 independent executions (3 per configuration, 30 rounds/run). Values of $\sigma_{\text{runs}}$ and $SE$ quantify dispersion and standard error across the 3 independent execution means ($SE = \sigma_{\text{runs}} / \sqrt{3}$).}\label{tab:tiempos_resumen}
\small
\setlength{\tabcolsep}{3pt}
\begin{tabular*}{\textwidth}{@{\extracolsep\fill}cccccccc}
\toprule
\textbf{Exp.} & \textbf{Nodes} & \textbf{Configuration} & \textbf{Runs} & \textbf{Median (s)} & \textbf{Mean (s)} & \textbf{$\sigma_{\text{runs}}$ (s)} & \textbf{SE (s)} \\
\midrule
$E_1$ & 1 & 4 val · with BC & 3 & \textbf{85.42} & 97.09 & 20.01 & 11.55 \\
$E_2$ & 1 & Baseline without BC & 3 & \textbf{75.70} & 76.45 & 1.65 & 0.95 \\
\midrule
$E_3$ & 2 & Baseline without BC & 3 & \textbf{66.62} & 67.62 & 2.29 & 1.32 \\
$E_4$ & 2 & 4 val · with BC & 3 & \textbf{72.96} & 73.29 & 1.28 & 0.74 \\
$E_5$ & 2 & 7 val · with BC & 3 & \textbf{73.35} & 74.40 & 0.54 & 0.31 \\
\bottomrule
\end{tabular*}
\end{table}

\begin{figure}[htbp]
    \centering
    \includegraphics[width=\columnwidth]{figuras/tiempo_ronda_todas.png}
    \caption{Comparison of client cycle time ($T_{\text{client}}$) across federated rounds for the five evaluated configurations.}
    \label{fig:tiempo_todas}
\end{figure}

The multi-run analysis reveals significant differences in operational stability across environments. In the local single-workstation configuration $E_1$, inter-run variability across execution means reached $\sigma_{\text{runs}} = 20.01$~s ($SE = 11.55$~s), with an aggregated longitudinal standard deviation across all 90 rounds of $\sigma_{\text{pool}} = 27.64$~s. This pronounced dispersion occurred because, during single-machine execution, local convolutional training competed directly for CPU and memory resources with eight background Docker containers hosting the Besu validator nodes and the IPFS daemon. Conversely, in the distributed deployment ($E_4$ and $E_5$), inter-run variability dropped markedly ($\sigma_{\text{runs}} \in [0.54, 1.28]$~s and $SE \le 0.74$~s; with pooled longitudinal standard deviations constrained to $\sigma_{\text{pool}} \le 3.92$~s), demonstrating excellent operational stability. Consequently, the \textbf{median} is adopted as the primary metric to isolate intrinsic system performance from transient external perturbations.

\subsection{Isolated Temporal Overhead of the Blockchain Layer}
To evaluate the precise temporal cost incurred by tensor serialization, IPFS upload, ECDSA signing, and IBFT 2.0 consensus, we compare corresponding configurations with and without the security layer (Figure~\ref{fig:coste_bc}):
\begin{itemize}
    \item \textbf{Single-Workstation Scenario ($E_1$ vs. $E_2$)}: Median cycle time increases from 75.70~s to 85.42~s, representing an overhead of \textbf{+9.72~s (+12.8\%)}.
    \item \textbf{Distributed Scenario ($E_3$ vs. $E_4$)}: Median cycle time increases from 66.62~s to 72.96~s, corresponding to an overhead of only \textbf{+6.34~s (+9.5\%)}.
\end{itemize}

\begin{figure}[htbp]
    \centering
    \includegraphics[width=\columnwidth]{figuras/coste_blockchain.png}
    \caption{Temporal overhead induced by blockchain integration in single-workstation and distributed deployments.}
    \label{fig:coste_bc}
\end{figure}

This finding is of particular significance for vehicular and edge systems: when offloading client training to dedicated hardware, the latency impact of the auditability layer remains strictly below 10\%. An additional cost of approximately 6 seconds per round represents a practical, viable trade-off in exchange for end-to-end immutability, non-repudiation, and verifiable integrity.

\subsection{Physical Distribution Effect and Phase Breakdown}
Comparing single-workstation and distributed configurations confirms that physical role decoupling significantly improves overall throughput (Figure~\ref{fig:coste_dist}): round latency decreases both in the baseline setup without blockchain (from 75.70~s to 66.62~s, a 12.0\% speedup) and with blockchain (from 85.42~s to 72.96~s, a 14.6\% speedup).

\begin{figure}[htbp]
    \centering
    \includegraphics[width=\columnwidth]{figuras/coste_distribucion.png}
    \caption{Effect of physical node distribution on cycle time: mitigating CPU/memory contention.}
    \label{fig:coste_dist}
\end{figure}

To analyze the underlying factors driving this behavior, Table~\ref{tab:desglose_fases} provides a fine-grained breakdown of average latency across individual phases at both clients (Station 2) and server (Station 1), demonstrating strict correspondence with the aggregate metrics in Table~\ref{tab:tiempos_resumen}.

\begin{table}[htbp]
\caption{Detailed breakdown of average phase latencies for clients ($T_{\text{client}}$) and coordinator ($T_{\text{server}}$) (in seconds).}\label{tab:desglose_fases}
\begin{tabular*}{\textwidth}{@{\extracolsep\fill}lcccc|cccc}
\toprule
& \multicolumn{4}{c|}{\textbf{Client Phases (Station 2)}} & \multicolumn{4}{c}{\textbf{Server Phases (Station 1)}} \\
\cmidrule(lr){2-5} \cmidrule(lr){6-9}
\textbf{Exp.} & \textbf{Train} & \textbf{Verif.} & \textbf{IPFS} & \textbf{Tx Besu} & \textbf{Wait} & \textbf{Verif.} & \textbf{FedAvg} & \textbf{Block} \\
\midrule
$E_1$ & 94.08 & 0.08 & 0.06 & 2.85 & 98.16 & 0.05 & 0.0065 & 5.14 \\
$E_2$ & 76.45 & 0.00 & 0.00 & 0.00 & 77.02 & 0.00 & 0.0040 & 0.00 \\
\midrule
$E_3$ & 67.62 & 0.00 & 0.00 & 0.00 & 71.07 & 0.00 & 0.0036 & 0.00 \\
$E_4$ & 66.22 & 1.68 & 2.43 & 2.84 & 79.54 & 0.06 & 0.0047 & 2.89 \\
$E_5$ & 65.40 & 1.91 & 3.93 & 2.94 & 82.63 & 0.06 & 0.0079 & 2.99 \\
\bottomrule
\end{tabular*}
\end{table}

The empirical breakdown in Table~\ref{tab:desglose_fases} substantiates three fundamental principles:
\begin{enumerate}
    \item \textbf{Computation decoupling:} In distributed configurations ($E_3$, $E_4$, $E_5$), time dedicated exclusively to PyTorch backpropagation remains constant ($\sim 65.4\text{--}67.6$~s), eliminating the CPU contention observed in $E_1$ ($94.08$~s).
    \item \textbf{Bounded network and verification overhead:} Remote JSON-RPC verification ($1.68\text{--}1.91$~s at clients, $\sim 5$~ms at coordinator), IPFS network transfer ($2.43\text{--}3.93$~s), and in-memory multihash calculation at the coordinator ($\sim 52$~ms in final code, $<0.08\%$ of the cycle) introduce minimal overhead, while FedAvg aggregation is practically instantaneous ($< 8$~ms).
    \item \textbf{Consistency between $T_{\text{client}}$ and $T_{\text{server}}$:} The sum of the four client-side local phases (\textit{Train} + \textit{Verif.} + \textit{IPFS} + \textit{Tx Besu}) strictly accounts for the client cycle time in Table~\ref{tab:tiempos_resumen} (e.g., in $E_4$: $66.22 + 1.68 + 2.43 + 2.84 = 73.17$~s, matching the empirical mean of $73.29$~s). Concurrently, server perspective reflects overall round duration ($T_{\text{server}}$), where the coordinator awaits concurrent client completion (\textit{Wait}, $79.54$~s) and commits the aggregated global model on-chain (\textit{Block}, $2.89$~s), totaling an end-to-end latency of $\approx 82.49$~s.
\end{enumerate}

\subsection{Impact of Validator Committee Size in IBFT 2.0 (4 vs. 7 Validators)}
In permissioned consortium deployments, elevating Byzantine fault tolerance requires expanding the validator committee such that $N \ge 3f + 1$, where $N$ is the total validator count and $f$ is the maximum number of tolerated faulty or adversarial nodes:
\begin{itemize}
    \item 4-validator network ($E_4$): tolerates $f=1$ faulty node; signature quorum $2f+1 = 3$.
    \item 7-validator network ($E_5$): tolerates $f=2$ faulty nodes; signature quorum $2f+1 = 5$.
\end{itemize}
Because all validator instances in our testbed execute containerized on the same physical host, this experiment does not aim to measure wide-area network scalability, but rather to isolate and quantify the direct impact of increasing the validator count and signature quorum on local consensus processing and round duration.

Comparing $E_4$ against $E_5$, median cycle time per round exhibits only a slight change, rising from 72.96~s to 73.35~s (+0.39~s), demonstrating that round cadence remains dominated by local model training rather than consensus messaging.

\begin{figure}[htbp]
    \centering
    \includegraphics[width=\columnwidth]{figuras/tiempo_confirmacion.png}
    \caption{Average on-chain transaction confirmation latency by topology and validator count in Hyperledger Besu.}
    \label{fig:tiempo_confirmacion}
\end{figure}

Isolating \textbf{on-chain transaction confirmation latency} (Figure~\ref{fig:tiempo_confirmacion}) reveals the empirical impact of expanding the consensus committee from 4 to 7 validator nodes in the evaluated testbed:
\begin{itemize}
    \item Client transaction (\texttt{registerClientModel}): mean confirmation latency increases from 2.84~s in $E_4$ to 2.94~s in $E_5$ (+0.10~s).
    \item Coordinator transaction (\texttt{updateModel}): mean confirmation latency increases from 2.89~s in $E_4$ to 2.99~s in $E_5$ (+0.10~s).
\end{itemize}
This symmetric 100~ms increase reflects the empirical consensus communication overhead observed when expanding the validator quorum ($f=1 \to 2$). This indicates that elevated Byzantine resilience can be readily accommodated without disrupting training throughput.

\subsection{Adversarial Robustness and Experimental Attack Verification}
\label{subsec:ataques}
To empirically verify that the architecture actively intercepts real-world adversarial vectors rather than merely logging nominal transactions, controlled attack scenarios were executed via the \texttt{test\_verificacion.py} test suite. The evaluation assesses system behavior against the two primary Byzantine threats in untrusted federated environments: coordinator equivocation or split-view attacks (downward flow) and unauthenticated or tampered client updates (upward flow).

\subsubsection{Test 1: Downward Verification (Tampered Server vs. Edge Client)}
\label{subsubsec:ataque_descendente}
In untrusted federated environments, distributing global weights exposes edge clients to two primary risks: in-transit parameter tampering by network intermediaries (\textit{Man-in-the-Middle}) and malicious server equivocation, where a compromised coordinator delivers divergent or deliberately poisoned models to specific participants. Crucially, the validity of the on-chain global model is strictly governed by the \texttt{onlyServer} access modifier in \texttt{FLRegistry}, which deterministically reverts unauthorized attempts to record global models. Consequently, the client-side downward verification routine autonomously safeguards edge nodes against discrepancies between gRPC-delivered weights and the authoritative blockchain commitment:
\begin{enumerate}
    \item \textbf{Generation and Authoritative Commitment:} The legitimate coordinator instantiates the \texttt{NetCliente} baseline model, serializes weights using \texttt{pickle.dumps}, and publishes them to the local IPFS daemon at \texttt{/api/v0/add?pin=false}, obtaining CID \texttt{QmQTpgejdzh...}. The coordinator then broadcasts a signed transaction calling \texttt{updateModel(round, official\_cid)} from its authorized address (\texttt{0xfe3b55...}), sealing the cryptographic commitment in \textbf{Block~80}.
    \item \textbf{Adversarial Injection and Discrepancy Simulation:} An attack scenario is simulated by injecting additive perturbation into parameters delivered to the client ($\mathbf{w}_{\text{cand}} = \mathbf{w}_{\text{official}} + 0.5$) and transmitting the perturbed vector via gRPC (modeling in-transit alteration or an equivocation attempt).
    \item \textbf{Cross-Verification and Immediate Abort:} Upon receiving the training request, the client queries \texttt{getLastModel()} on the smart contract via JSON-RPC to retrieve the on-chain committed CID. It concurrently downloads the reference binary artifact from the IPFS gateway (\texttt{/ipfs/<CID>}) and performs an element-wise numerical tensor comparison within floating-point tolerance (\texttt{np.allclose(p\_received, p\_ipfs, atol=1e-7)}).
\end{enumerate}

\begin{center}
\small
\fbox{
\begin{minipage}{0.95\columnwidth}
\textbf{Algorithm 1. Downward Verification at the Edge Client}
\begin{enumerate}
    \item \textbf{Input:} Candidate weights $\mathbf{w}_{\text{cand}}$ received via gRPC from coordinator.
    \item Query Besu smart contract: $\text{CID}_{\text{onchain}} \leftarrow \texttt{FLRegistry.getLastModel().cid}$
    \item Fetch binary model from IPFS: $\mathbf{w}_{\text{ipfs}} \leftarrow \texttt{IPFS.get}(\text{CID}_{\text{onchain}})$
    \item \textbf{If} $\exists l : \|\mathbf{w}_{\text{cand}, l} - \mathbf{w}_{\text{ipfs}, l}\|_{\infty} > \epsilon$ \textbf{then}
    \begin{itemize}
        \item Emit alert: \textit{``Discrepancy detected in received weights''}
        \item \textbf{ABORT} local round immediately (block backpropagation)
    \end{itemize}
    \item \textbf{Else:} Accept $\mathbf{w}_{\text{cand}}$ and proceed with local training.
\end{enumerate}
\end{minipage}
}
\end{center}

Figure~\ref{fig:log_ataque_descendente} shows runtime console telemetry at the edge client. The complete verification sequence (RPC query to Besu, IPFS download, and tensor comparison) required only \textbf{0.1486~seconds} (149~ms). Upon detecting the discrepancy, the client immediately terminated execution, preventing local model corruption and safeguarding private vehicular training data.

\begin{figure}[htbp]
    \centering
    \includegraphics[width=\columnwidth]{figuras/logs_deteccion_ataque_descendent.png}
    \caption{Test 1: Console log of edge client detecting parameter tampering against the official on-chain CID sealed in block 80.}
    \label{fig:log_ataque_descendente}
\end{figure}

Figure~\ref{fig:bloque_80_audit} shows the immutable blockchain evidence captured from the Hyperledger Besu block explorer. It documents the confirmation of transaction \texttt{updateModel} in \textbf{Block~80}, submitted by coordinator address \texttt{0xfe3b557e...} with transaction hash \texttt{0x973b0b...} and gas consumption of 167,460 in this preliminary run (which becomes 169,649 in the final contract featuring \texttt{onlyServer}, gas limit 500,000). This bijective correspondence between the mined block and client verification confirms protocol consistency.

\begin{figure}[htbp]
    \centering
    \includegraphics[width=\columnwidth]{figuras/auditoria_grafana_bloque_80.png}
    \caption{Test 1: Confirmation of Block 80 in Hyperledger Besu, recording the authoritative \texttt{updateModel} transaction submitted by the coordinator.}
    \label{fig:bloque_80_audit}
\end{figure}

\subsubsection{Test 2: Upward Verification (Tampered Client Update vs. Aggregation Coordinator)}
\label{subsubsec:ataque_ascendente}
The second critical vulnerability is the injection of manipulated, tampered, or unregistered parameter weights by a client that do not match its registered on-chain commitment. To intercept this at the gateway prior to FedAvg aggregation, the coordinator cross-verifies the bijective correspondence between received memory payloads and signed on-chain commitments:
\begin{enumerate}
    \item \textbf{Legitimate On-Chain Registration:} The honest client node (account 3, address \texttt{0xf17f52...}) uploads its local weights to IPFS and submits a signed transaction calling \texttt{registerClientModel(round, legitimate\_cid)}, confirmed in \textbf{Block~81}.
    \item \textbf{Adversarial Injection:} An injection attempt is simulated where an update with altered parameters or mismatched metadata is delivered to the coordinator (modeling the transmission of a fraudulent CID \texttt{QmClientFraud...} or tampered weight tensors that deviate from the registered on-chain digest).
    \item \textbf{Cryptographic Audit and Exclusion from FedAvg:} Prior to executing FedAvg, the coordinator computes the cryptographic hash of received memory tensors ($\text{CID}_{\text{calc}}^{(k)}$) and executes a read-only JSON-RPC call \texttt{getClientModel(round, address)} against \texttt{FLRegistry} to retrieve $\text{CID}_{\text{onchain}}^{(k)}$. Upon detecting discrepancy, the update is excluded before aggregation.
\end{enumerate}

\begin{center}
\small
\fbox{
\begin{minipage}{0.95\columnwidth}
\textbf{Algorithm 2. Upward Cryptographic Audit and RPC Check at Coordinator}
\begin{enumerate}
    \item \textbf{Input:} Weights $\mathbf{w}_k^{(t)}$ and identifier $\text{CID}_{\text{rep}}^{(k)}$ reported by client $k$ with address $\text{addr}_k$.
    \item Compute hash digest of received in-memory weights: \\
    $\text{CID}_{\text{calc}}^{(k)} \leftarrow \text{Hash}(\mathbf{w}_k^{(t)})$
    \item Query client on-chain record from Besu: \\
    $\text{CID}_{\text{onchain}}^{(k)} \leftarrow \texttt{getClientModel}(t, \text{addr}_k)$
    \item \textbf{If} $\text{CID}_{\text{rep}}^{(k)} \neq \text{CID}_{\text{onchain}}^{(k)}$ \textbf{or} $\text{CID}_{\text{calc}}^{(k)} \neq \text{CID}_{\text{onchain}}^{(k)}$ \textbf{then}
    \begin{itemize}
        \item Emit alert: \textit{``Discrepancy detected / Unauthenticated injection''}
        \item \textbf{DISCARD} client $k$ contribution from FedAvg pool
    \end{itemize}
    \item \textbf{Else:} Include weights $\mathbf{w}_k^{(t)}$ in global aggregation.
\end{enumerate}
\end{minipage}
}
\end{center}

Figure~\ref{fig:log_ataque_ascendente} illustrates server terminal output. The JSON-RPC verification call to the EVM required only \textbf{0.0309~seconds} (31~ms). Confirming that the submitted CID did not match the client signed transaction, the server discarded the rogue update immediately, preserving global model integrity.

\begin{figure}[htbp]
    \centering
    \includegraphics[width=\columnwidth]{figuras/logs_deteccion_ataque_ascendente.png}
    \caption{Test 2: Coordinator console log auditing client CID via RPC and discarding fraudulent contribution prior to aggregation.}
    \label{fig:log_ataque_ascendente}
\end{figure}

Figure~\ref{fig:bloque_81_audit} displays the transaction receipt in \textbf{Block~81} from the Besu explorer, documenting the transaction \texttt{registerClientModel} issued by client account \texttt{0xf17f5215...} with hash \texttt{0x347899f...} and gas consumption of 160,200 (gas limit 200,000). Because the adversary transmitted a CID divergent from that sealed on-chain, the cross-check intercepted the injection due to the collision resistance of the underlying SHA-256 multihash.

\begin{figure}[htbp]
    \centering
    \includegraphics[width=\columnwidth]{figuras/auditoria_grafana_bloque_81.png}
    \caption{Test 2: Forensic confirmation in Hyperledger Besu of Block 81, recording legitimate client transaction \texttt{registerClientModel}.}
    \label{fig:bloque_81_audit}
\end{figure}

\subsubsection{Test 3: Ex-Post Auditing and Independent FedAvg Verifiability}
\label{subsubsec:auditoria_expost_eval}
To evaluate the third integrity property formalised in Section~\ref{subsec:auditoria_fedavg}, we assessed the capability of an external auditor or regulator to retrospectively (\textit{ex-post}) and autonomously verify the central server aggregation fidelity. This test verifies whether the coordinator computed FedAvg honestly without injecting unverified bias, omitting clients, or altering weights.

The audit was executed on Round~2 of a federated training run comprising two concurrent edge clients connected to Hyperledger Besu (contract \texttt{0x42699A76...}) and IPFS. Table~\ref{tab:trazabilidad_ronda2} summarizes the complete chain of custody for Round~2, correlating sending accounts, transactions confirmed across consecutive blocks 73, 74, and 75, IPFS multihash identifiers, and execution latencies.

\begin{table}[htbp]
\caption{Chain of custody and cross-traceability on-chain/IPFS during Round 2 and ex-post audit.}\label{tab:trazabilidad_ronda2}
\footnotesize
\setlength{\tabcolsep}{3pt}
\renewcommand{\arraystretch}{1.2}
\begin{tabular*}{\textwidth}{@{\extracolsep\fill}lcccccc}
\toprule
\textbf{Entity} & \makecell{\textbf{Ethereum}\\\textbf{Address}} & \makecell{\textbf{Event /}\\\textbf{Method}} & \makecell{\textbf{Besu}\\\textbf{Block}} & \textbf{Latency} & \makecell{\textbf{IPFS}\\\textbf{CID}} & \makecell{\textbf{State /}\\\textbf{Result}} \\
\midrule
Client 3 & \texttt{0xf17f5215...} & \makecell[c]{\texttt{registerClient-}\\\texttt{Model}} & Block 73 & 0.886 s & \texttt{QmS5ViVf...} & \makecell[c]{Confirmed\\on-chain} \\
Client 2 & \texttt{0x62730609...} & \makecell[c]{\texttt{registerClient-}\\\texttt{Model}} & Block 74 & 4.311 s & \texttt{QmNh3TLS...} & \makecell[c]{Confirmed\\on-chain} \\
Server  & \texttt{0xfe3b557e...} & \makecell[c]{\texttt{updateModel}\\(FedAvg)} & Block 75 & 4.709 s & \texttt{QmWF7tJ1...} & \makecell[c]{Global model\\committed} \\
\midrule
Auditor & \textit{External (Off-chain)} & \makecell[c]{FedAvg\\Reconstruction} & \makecell[c]{Query\\73--75} & 0.339 s & \texttt{QmWF7tJ1...} & \makecell[c]{\textbf{$\Delta = 0.000000$}\\\textbf{(MATCH)}} \\
\bottomrule
\end{tabular*}
\end{table}

The audit protocol proceeded through the following steps:
\begin{enumerate}
    \item \textbf{On-chain extraction and indexing:} In an off-band, decoupled manner, the auditor dynamically identifies participating clients $\mathcal{K}_2$ by querying \texttt{ModelUpdated} event logs for Round~2 in Besu (retrieving Client~2 and Client~3 accounts) and extracts their registered CIDs (\texttt{QmNh3TLS...} and \texttt{QmS5ViVf...}), along with the coordinator global model commitment (\texttt{QmWF7tJ1...}).
    \item \textbf{Deterministic retrieval and reconstruction:} The auditor downloads all three binary model artifacts from IPFS ($2.08$~MB per model, totaling $6.24$~MB retrieved in an accumulated $0.200$~s [$0.096$~s, $0.044$~s, and $0.060$~s]) and executes canonical FedAvg aggregation:
    \begin{equation}
    \mathbf{w}_{\text{recon}}^{(2)} = \frac{1}{2}\mathbf{w}_{\text{cli2}}^{(2)} + \frac{1}{2}\mathbf{w}_{\text{cli3}}^{(2)}
    \end{equation}
    generating the resulting model CID using the SHA-256 multihash algorithm.
    \item \textbf{Tensor and cryptographic verification:} The reconstructed parameter tensors are compared layer-by-layer against the global model sealed by the coordinator in Block~75.
\end{enumerate}

\begin{figure}[htbp]
    \centering
    \includegraphics[width=\columnwidth]{figuras/logs_auditoria_expost_fedavg.png}
    \caption{Test 3: Independent auditor console log reconstructing Round 2 FedAvg aggregation from Besu and IPFS, verifying exact fidelity ($\Delta = 0.000000$) and CID match.}
    \label{fig:log_auditoria_expost}
\end{figure}

\begin{sidewaystable}
\caption{Experimental validation of the integrity and auditing framework across manipulation and ex-post verification scenarios.}\label{tab:resumen_ataques}
\footnotesize
\setlength{\tabcolsep}{6pt}
\renewcommand{\arraystretch}{1.25}
\begin{tabular*}{\textheight}{@{\extracolsep\fill}lcccc}
\toprule
\textbf{Test Scenario} & \makecell{\textbf{Defending / Auditing}\\\textbf{Entity}} & \makecell{\textbf{Preventive Action /}\\\textbf{Result}} & \textbf{Latency} & \makecell{\textbf{Weight}\\\textbf{Integrity}} \\
\midrule
\makecell[l]{In-Transit Tampering /\\Untrusted Coordinator} & \makecell{Client\\(Edge)} & \makecell{Deterministic interception /\\Preventive abort} & 0.1486 s & \makecell{No local\\contamination} \\ \midrule
\makecell[l]{Unauthenticated Injection /\\Unregistered Client} & \makecell{Aggregation\\Coordinator} & \makecell{Deterministic interception /\\FedAvg exclusion} & 0.0309 s & \makecell{Aggregation\\preserved} \\ \midrule
\makecell[l]{Ex-Post FedAvg\\Audit} & \makecell{Independent\\Auditor} & \makecell{Mathematical reconstruction /\\Faithful verification} & 0.34 s & \makecell{Exact fidelity\\($\Delta = 0$, CID MATCH)} \\
\bottomrule
\end{tabular*}
\end{sidewaystable}

As documented in the console telemetry in Figure~\ref{fig:log_auditoria_expost}, the entire audit sequence (dynamic participant discovery, RPC query of commitments, downloading 6.24~MB from IPFS, tensor reconstruction, and verification across all 8 layers of the CNN) required only \textbf{0.3391~seconds}. Parametric discrepancy across all layers was strictly zero ($\Delta = 0.000000\text{e}+00$), resulting in an identical CID to that registered in Block~75 (\texttt{CID Equality: MATCH}), confirming aggregation correctness.

Furthermore, mechanism sensitivity was evaluated under simulated adversarial mode (\texttt{--simulate-tampering}). By injecting an additive perturbation of $+0.1$ into coordinator weights, the auditor detected the anomaly in \textbf{0.1805~seconds} ($\Delta = 1.000000\text{e}-01$), logging a formal discrepancy alert (\texttt{AUDIT FAILED: DISCREPANCY DETECTED}).

\subsubsection{Synthesis of Validation and Security Scope}
Table~\ref{tab:resumen_ataques} synthesizes the quantitative metrics obtained across the three experimental integrity and auditing tests.

These experiments provide functional validation of the cryptographic primitives: when encountering altered in-transit parameters or uncommitted models, the protocol intercepts anomalies deterministically with response times in the order of tenths of a second (149~ms at the client, 31~ms at the coordinator, and $0.34$~s in ex-post FedAvg auditing), preventing corrupt weights from contaminating training.

As delineated in our threat model (Section~\ref{subsec:modelo_amenazas}), operational boundaries must be emphasized: on-chain verification counters parameter tampering in transit, unauthorized submissions via ECDSA signatures, and rogue uncommitted models. However, if a malicious node generates poisoned updates from local training data (\textit{data poisoning} or \textit{label flipping})~\cite{tolpegin2020data} and legitimately signs the transaction in Besu, cryptographic commitments will match despite semantic contamination. Mitigating semantic attacks requires augmenting traceability with Byzantine-robust aggregation rules (e.g., Krum~\cite{blanchard2017byzantine}, Median, or Trimmed-Mean) or conducting forensic ex-post auditing enabled by our architecture (Section~\ref{subsec:auditoria_fedavg}).

\subsection{EVM Gas Consumption and Predictability}
Figure~\ref{fig:consumo_gas} illustrates EVM gas consumption for operations recorded across rounds.

\begin{figure}[htbp]
    \centering
    \includegraphics[width=0.70\columnwidth]{figuras/consumo_gas.png}
    \caption{Ethereum Virtual Machine gas consumption by transaction type across 30 federated rounds.}
    \label{fig:consumo_gas}
\end{figure}

The gas profile is invariant across all training iterations:
\begin{itemize}
    \item \texttt{registerClientModel}: exactly \textbf{160,200 gas units}.
    \item \texttt{updateModel}: exactly \textbf{169,649 gas units}.
\end{itemize}
This consumption reflects the final version of the \texttt{FLRegistry} smart contract: incorporating the \texttt{onlyServer} access modifier into \texttt{updateModel} introduces an EVM identity check (\texttt{SLOAD} and conditional branch on \texttt{serverAddress}) that establishes execution cost at exactly 169,649 gas (a bounded overhead of 2,189 gas, +1.3\% compared to an open contract without access control), remaining strictly constant and predictable across all operational rounds. The sole observed exception occurs in Round~0 for the coordinator, which required \textbf{186,749 gas units}. This discrepancy (+17,100 gas) corresponds to the EVM \texttt{SSTORE} instruction when allocating uninitialized storage slots (\textit{cold write}) and toggling the \texttt{hasModel} boolean flag.

This consumption invariance formally validates the decoupled architecture: by persisting only uniform-length strings on-chain (46-character IPFS CIDv0 identifiers encapsulating SHA-256 multihashes and bounded metadata), gas costs remain constant per transaction and independent of neural network parameter count and cumulative round count, guaranteeing predictable resource utilization in consortium networks.

\subsection{On-Chain Immutability Audit and Forensic Traceability}
\label{subsec:auditoria_onchain}
To empirically verify bijective correspondence between Flower lifecycle events and immutable blockchain persistence, Table~\ref{tab:auditoria_transacciones} presents a representative sample of transaction receipts extracted from runtime telemetry, correlated directly with the block explorer visualization (Figure~\ref{fig:auditoria_bloque}).

\begin{table}[htbp]
\caption{Sample of on-chain forensic traceability and audit receipts.}\label{tab:auditoria_transacciones}
\footnotesize
\setlength{\tabcolsep}{4pt}
\renewcommand{\arraystretch}{1.15}
\begin{tabular*}{\textwidth}{@{\extracolsep\fill}cccccc}
\toprule
\textbf{Block} & \textbf{Entity} & \textbf{Event / Method} & \textbf{Sender Address} & \textbf{Tx Hash Prefix} & \textbf{EVM Gas} \\
\midrule
8 & Coordinator & Deploy \texttt{FLRegistry} & \texttt{0xfe3b557e...} & \texttt{0x939eb46c...} & 1,999,000 \\
13 & Coordinator & \texttt{updateModel} & \texttt{0xfe3b557e...} & \texttt{0xa18e3bb4...} & 169,649 \\
14 & Client 2 & \texttt{registerClientModel} & \texttt{0x62730609...} & \texttt{0xc1c7402f...} & 160,200 \\
\midrule
765 & Coordinator & \texttt{updateModel} (Round 23) & \texttt{0xfe3b557e...} & \texttt{0x4ceabb7a...} & 169,649 \\
783 & Coordinator & \texttt{updateModel} (Round 24) & \texttt{0xfe3b557e...} & \texttt{0xb9cb75f5...} & 169,649 \\
801 & Coordinator & \texttt{updateModel} (Round 25) & \texttt{0xfe3b557e...} & \texttt{0xdf69295f...} & 169,649 \\
\bottomrule
\end{tabular*}
\end{table}

\begin{figure}[htbp]
    \centering
    \includegraphics[width=0.96\columnwidth]{figuras/bloque_contrato.png}
    \caption{On-chain immutability audit: Besu explorer visualization of \texttt{FLRegistry} contract deployment (Block 8) by coordinator (\texttt{0xfe3b55...}) alongside operational transactions 13 and 14.}
    \label{fig:auditoria_bloque}
\end{figure}

Each block identifier and cryptographic transaction hash provides non-repudiable forensic proof, confirming the viability of the architecture for regulatory audits in vehicular and safety-critical distributed systems.
